\subsection{Outer Product of Kets}

In quantum mechanics, \definition{Bra-Ket notation} (introduced by Dirac) is a concise and powerful tool for representing quantum states and operations. It uses \textbf{kets} \ket{\psi} for \textbf{vectors} (states), and \textbf{bras} \bra{\psi} for their \textbf{Hermitian adjoints} (dual vectors).

\highspace
Given two kets:
\begin{equation*}
    \ket{a} = \begin{pmatrix}
        a_{1} \\
        a_{2}
    \end{pmatrix}
    \hspace{2em}
    \ket{b} = \begin{pmatrix}
        b_{1} \\
        b_{2}
    \end{pmatrix}
\end{equation*}
The \definitionWithSpecificIndex{Outer Product}{Outer Product of Kets}{} \ket{a}\bra{b} is a \textbf{matrix}, formed by:
\begin{equation}
    \ket{a}\bra{b} = \ket{a} \otimes \bra{b} = \begin{pmatrix}
        a_{1} \\
        a_{2}
    \end{pmatrix}
    \begin{pmatrix}
        \overline{b_{1}} & \overline{b_{2}}
    \end{pmatrix}
    =
    \begin{pmatrix}
        a_{1}\overline{b_{1}} & a_{1}\overline{b_{2}} \\
        a_{2}\overline{b_{1}} & a_{2}\overline{b_{2}}
    \end{pmatrix}
\end{equation}
Where the symbol $\otimes$ is the \textbf{tensor product operator}. This operation \textbf{maps a vector to a matrix}.

\highspace
\textcolor{Green3}{\faIcon{question-circle} \textbf{Why is the outer product useful?}} Because \textbf{all quantum gates can be written using outer products of kets}. For example, the \textbf{Pauli-X gate} can be expressed as:
\begin{equation*}
    X = \ket{0}\bra{1} + \ket{1}\bra{0} =
    \underbrace{\begin{pmatrix}
        1 \\ 0
    \end{pmatrix}
    \begin{pmatrix}
        0 & 1
    \end{pmatrix}}_{\begin{pmatrix}
        0 & 1 \\
        0 & 0
    \end{pmatrix}}
    +
    \underbrace{\begin{pmatrix}
        0 \\ 1
    \end{pmatrix}
    \begin{pmatrix}
        1 & 0
    \end{pmatrix}}_{\begin{pmatrix}
        0 & 0 \\
        1 & 0
    \end{pmatrix}}
    =
    \begin{pmatrix}
        0 & 1 \\
        1 & 0
    \end{pmatrix}
\end{equation*}

\highspace
The \definitionWithSpecificIndex{Inner Product}{Inner Product of Kets}{} \braket{b}{a} is a scalar:
\begin{equation}
    \braket{b}{a} = \overline{b_{1}} a_{1} + \overline{b_{2}} a_{2}
\end{equation}
So:
\begin{itemize}
    \item \textbf{Inner product} $\rightarrow$ \textbf{projection}, result is number (scalar).
    \item \textbf{Outer product} $\rightarrow$ \textbf{matrix}, result is operator.
\end{itemize}

\highspace
\begin{flushleft}
    \textcolor{Green3}{\faIcon{book} \textbf{Associativity between outer and inner products}}
\end{flushleft}
A key property of the outer product is how it interacts with inner products. Specifically:
\begin{equation}
    \left(\ket{a}\bra{b}\right) \ket{c} = \ket{a} \left( \braket{b}{c} \right)
\end{equation}
This identity shows that the outer product acts as an operator: it first computes the inner product $\braket{b}{c}$ (a scalar), and then scales the vector $\ket{a}$.

\begin{proof}[\textbf{Proof the Associativity}]
    In this proof, we will demonstrate that applying the outer product $\ket{a}\bra{b}$ to a ket $\ket{c}$ is equivalent to scaling $\ket{a}$ by the inner product $\braket{b}{c}$.

    \highspace
    Let's define:
    \begin{equation*}
        \ket{a} = \begin{pmatrix}
            a_{1} \\
            a_{2}
        \end{pmatrix}
        \hspace{2em}
        \ket{b} = \begin{pmatrix}
            b_{1} \\
            b_{2}
        \end{pmatrix}
        \hspace{2em}
        \ket{c} = \begin{pmatrix}
            c_{1} \\
            c_{2}
        \end{pmatrix}
    \end{equation*}
    \begin{enumerate}
        \item Compute \textbf{inner product} \braket{b}{c}:
        \begin{equation*}
            \braket{b}{c} = \overline{b_{1}} c_{1} + \overline{b_{2}} c_{2} = \alpha
        \end{equation*}
        Where $\alpha$ is a \textbf{scalar} value.

        \item Multiply \textbf{scalar} $\alpha$ with \ket{a} (vector):
        \begin{equation*}
            \ket{a}\alpha = \begin{pmatrix}
                a_{1} \\
                a_{2}
            \end{pmatrix}\alpha
            =
            \begin{pmatrix}
                a_{1} \cdot \alpha \\
                a_{2} \cdot \alpha
            \end{pmatrix}
            =
            \begin{pmatrix}
                a_{1} \cdot \left(\overline{b_{1}} c_{1} + \overline{b_{2}} c_{2}\right) \\
                a_{2} \cdot \left(\overline{b_{1}} c_{1} + \overline{b_{2}} c_{2}\right)
            \end{pmatrix}
        \end{equation*}

        \item Compute outer product \ket{a}\bra{b} (matrix):
        \begin{equation*}
            \ket{a}\bra{b} =
            \begin{pmatrix}
                a_{1} \\
                a_{2}
            \end{pmatrix}
            \begin{pmatrix}
                \overline{b_{1}} & \overline{b_{2}}
            \end{pmatrix}
            =
            \begin{pmatrix}
                a_{1}\overline{b_{1}} & a_{1}\overline{b_{2}} \\
                a_{2}\overline{b_{1}} & a_{2}\overline{b_{2}}
            \end{pmatrix}
        \end{equation*}

        \item Multiply \ket{a}\bra{b} by \ket{c}:
        \begin{equation*}
            \left(\ket{a}\bra{b}\right) \ket{c} =
            \begin{pmatrix}
                a_{1}\overline{b_{1}} & a_{1}\overline{b_{2}} \\
                a_{2}\overline{b_{1}} & a_{2}\overline{b_{2}}
            \end{pmatrix}
            \begin{pmatrix}
                c_{1} \\
                c_{2}
            \end{pmatrix}
            =
            \begin{pmatrix}
                a_{1} \cdot \left( \overline{b_{1}} c_{1} + \overline{b_{2}} c_{2} \right) \\
                a_{2} \cdot \left( \overline{b_{1}} c_{1} + \overline{b_{2}} c_{2} \right)
            \end{pmatrix}
        \end{equation*}
    \end{enumerate}
    The second and fourth steps yield the same result.
    This proves that \textbf{outer product applied to a ket is equivalent to inner product \braket{b}{c} scaled by \ket{a}}.
\end{proof}
